\scanpage{73}
July 18, I962\\

My dear Murre,\\

I recently had some thought on finiteness conditions for Picard preschemes, and substantially improved on the results stated in the last section of my last Bourbaki talk. The main result stated there for a simple projective morphism with connected geometric fibers (namely that the pieces $\underline{\mathrm{pic}}^{P}_{X/S}$ are of finite type over $S$) has been extended by Mumford to the case where instead of $f$ simple we assume only $f$ flat with integral geometric fibers,(at least if these are normal). Using his result (the proof of which is quite simple and beautiful) , I could get rid of the normality assumption, and even (as in theorem 4.1. of my talk) restrict to the consideration of the two first non t\unsure{ri|\deleted{r}i|\deleted{i}r}vial coefficients of the Hilbert polynomials. The key results for the reduction are the following (the proofs being very technical, and rather different for (\unsure{\deleted{a}i|\deleted{i}a|i}) and (ii), except that (ii) uses (i) to reduce to the normal case; moreover (ii) uses Mumford's result and the equivalence criteria as developped in my last Seminar):

(i) Let $X$, $Y$ be proper over $S$ neetherian, let $f:X\handwritten{\to}Y$ be a \uline{surjective} $S$ - morphism, assume for simplicity of the statement that the Picard preschemes exist, the\handwritten{n} $f:\underline{\mathrm{Pic}}_{Y/S}\handwritten{\longrightarrow}\underline{\mathrm{Pic}}_{X/S}$ is of finite type (and in fact affine if $S$ is the spectrum of a field), i. e.a subset $M$ of $\underline{\mathrm{Pic}}_{Y/S}$ is quasi-compact iff its image in $\underline{\mathrm{Pic}}_{X/S}$ is.

(ii) The same conclusion holds for a canonical immersion $X\handwritten{\to}Y$, if $Y/S$ is projective with fibers all components of which are of dimension $\handwritten{\geqslant}3$, and if $X$ is the sub-scheme of zeros of a section over $Y$ of an invertible sheaf $\underline{L}$ ample with respect to $S$.

A connected result is that for any $X/S$ proper, and integer $n\handwritten{\ne}0$, the $n$. th power homomorphism in the Picard prescheme is of finite type.

I tell you about this, namely (i), because of the method of proof, involving of course considerations of non flat descent. The fact that I do not have any good effectivity criterion does not hamper, by just recalling what the effectivity of a given descent datum means. Now it turns out that by a slightly more careful analysis of the situation, one can prove the following theorem, of a type very close to the one you have proved recently, and to some you still want to prove as I understand it.

\hfill .../...\par
% archive sheet number: 72
